\documentclass[reqno]{amsart}
\usepackage[T1]{fontenc}
\usepackage{amsmath,amssymb,amsthm}
\usepackage{xr-hyper}
\usepackage[hidelinks]{hyperref}
\externaldocument[ip-]{20261124T214600--the-isotropic-plane__workbook}
\externaldocument[ib-]{20261125T001100--inert-primes-and-estimate-b__workbook}

% Chapter number; \ref{chap} gives it.
\newcounter{chap}
% Chapter numbers, generated from 20261125T012000--grh-workbook__guide.org
% by scripts/make-chapters.py.  Do not edit: change the order of the
% index and rerun the script (./preview-site.sh and CI do this).
%
% Usage in a sheet, after \newcounter{chap}:
%   % Chapter numbers, generated from 20261125T012000--grh-workbook__guide.org
% by scripts/make-chapters.py.  Do not edit: change the order of the
% index and rerun the script (./preview-site.sh and CI do this).
%
% Usage in a sheet, after \newcounter{chap}:
%   % Chapter numbers, generated from 20261125T012000--grh-workbook__guide.org
% by scripts/make-chapters.py.  Do not edit: change the order of the
% index and rerun the script (./preview-site.sh and CI do this).
%
% Usage in a sheet, after \newcounter{chap}:
%   \input{chapters}\chapnumber{<sheet ID>}
\makeatletter
\providecommand\chapters@set[2]{\@namedef{chapters@#1}{#2}}
\chapters@set{20261125T001000}{1}% zeros-to-ramp-sums
\chapters@set{20261125T001100}{2}% inert-primes-and-estimate-b
\chapters@set{20261125T005000}{3}% galois-fourier-and-averaging
\chapters@set{20261125T005500}{4}% columns-from-a-sphere
\chapters@set{20261125T000000}{5}% slope-method
\chapters@set{20261126T030100}{6}% a-twist-for-each-prime
\chapters@set{20261125T003000}{7}% the-chain-and-the-square-root
\chapters@set{20261125T010000}{8}% rank-lemma-and-zero-estimate
\chapters@set{20261125T010500}{9}% poincare-series-and-near-orthogonality
\chapters@set{20261124T223000}{10}% elliptic-curves-for-the-proof
\chapters@set{20261125T009500}{11}% fock-model-and-theta-intertwiners
\chapters@set{20261124T223100}{12}% complex-multiplication-and-inert-reduction
\chapters@set{20261125T008000}{13}% one-cm-determinant
\chapters@set{20261124T214500}{14}% supersingular-theta-divisibility
\chapters@set{20261124T214600}{15}% the-isotropic-plane
\chapters@set{20261125T013700}{16}% pairings-on-a-formal-group-of-height-two
\chapters@set{20261125T013500}{17}% hecke-nondegeneracy-on-the-tree
\chapters@set{20261126T030200}{18}% the-bound-is-attained
\chapters@set{20261125T009000}{19}% kronecker-limit-and-chowla-selberg
\chapters@set{20261125T002000}{20}% definite-theta-lifts
\chapters@set{20261125T015500}{21}% siegel-weil-in-the-smallest-case
\chapters@set{20261125T002100}{22}% definite-theta-lifts-and-norms
\chapters@set{20261125T016500}{23}% hecke-operators-and-theta-lifts
\chapters@set{20261125T016700}{24}% divisibility-along-the-chain
\chapters@set{20261125T004000}{25}% denominators-and-sha
\chapters@set{20261125T004100}{26}% tate-twists-and-sha
\chapters@set{20261126T020430}{27}% the-minus-part-by-gauss-sums
\chapters@set{20261125T014100}{28}% congruence-modules
\chapters@set{20261125T014500}{29}% eisenstein-congruences-and-selmer-classes
\chapters@set{20261125T014700}{30}% pseudocharacters-and-the-length-bound
\chapters@set{20261125T011000}{31}% galois-representations-black-box
\chapters@set{20261126T030300}{32}% the-data
\chapters@set{20261126T020210}{33}% bessel-poincare-series-and-their-gram-operators
\chapters@set{20261126T030400}{34}% the-seesaw-identity
\chapters@set{20261126T030500}{35}% divisibility-of-the-composite-lift
\chapters@set{20261126T030600}{36}% theorem-a-short-forms
\chapters@set{20261126T030700}{37}% the-determinant-of-a-chain
\chapters@set{20261126T030800}{38}% the-cost-of-one-lift
\chapters@set{20261126T030900}{39}% denominators-of-projections-to-lifts
\chapters@set{20261126T031000}{40}% twists-with-few-congruences
\chapters@set{20261126T031100}{41}% theorem-c-and-the-main-inequality
\providecommand\chapnumber[1]{%
  \@ifundefined{chapters@#1}%
    {\PackageError{chapters}{Sheet #1 is not listed in the index}%
      {List it in 20261125T012000--grh-workbook__guide.org\MessageBreak
       and run scripts/make-chapters.py.}}%
    {\edef\thechap{\@nameuse{chapters@#1}}}}
\makeatother
\chapnumber{<sheet ID>}
\makeatletter
\providecommand\chapters@set[2]{\@namedef{chapters@#1}{#2}}
\chapters@set{20261125T001000}{1}% zeros-to-ramp-sums
\chapters@set{20261125T001100}{2}% inert-primes-and-estimate-b
\chapters@set{20261125T005000}{3}% galois-fourier-and-averaging
\chapters@set{20261125T005500}{4}% columns-from-a-sphere
\chapters@set{20261125T000000}{5}% slope-method
\chapters@set{20261126T030100}{6}% a-twist-for-each-prime
\chapters@set{20261125T003000}{7}% the-chain-and-the-square-root
\chapters@set{20261125T010000}{8}% rank-lemma-and-zero-estimate
\chapters@set{20261125T010500}{9}% poincare-series-and-near-orthogonality
\chapters@set{20261124T223000}{10}% elliptic-curves-for-the-proof
\chapters@set{20261125T009500}{11}% fock-model-and-theta-intertwiners
\chapters@set{20261124T223100}{12}% complex-multiplication-and-inert-reduction
\chapters@set{20261125T008000}{13}% one-cm-determinant
\chapters@set{20261124T214500}{14}% supersingular-theta-divisibility
\chapters@set{20261124T214600}{15}% the-isotropic-plane
\chapters@set{20261125T013700}{16}% pairings-on-a-formal-group-of-height-two
\chapters@set{20261125T013500}{17}% hecke-nondegeneracy-on-the-tree
\chapters@set{20261126T030200}{18}% the-bound-is-attained
\chapters@set{20261125T009000}{19}% kronecker-limit-and-chowla-selberg
\chapters@set{20261125T002000}{20}% definite-theta-lifts
\chapters@set{20261125T015500}{21}% siegel-weil-in-the-smallest-case
\chapters@set{20261125T002100}{22}% definite-theta-lifts-and-norms
\chapters@set{20261125T016500}{23}% hecke-operators-and-theta-lifts
\chapters@set{20261125T016700}{24}% divisibility-along-the-chain
\chapters@set{20261125T004000}{25}% denominators-and-sha
\chapters@set{20261125T004100}{26}% tate-twists-and-sha
\chapters@set{20261126T020430}{27}% the-minus-part-by-gauss-sums
\chapters@set{20261125T014100}{28}% congruence-modules
\chapters@set{20261125T014500}{29}% eisenstein-congruences-and-selmer-classes
\chapters@set{20261125T014700}{30}% pseudocharacters-and-the-length-bound
\chapters@set{20261125T011000}{31}% galois-representations-black-box
\chapters@set{20261126T030300}{32}% the-data
\chapters@set{20261126T020210}{33}% bessel-poincare-series-and-their-gram-operators
\chapters@set{20261126T030400}{34}% the-seesaw-identity
\chapters@set{20261126T030500}{35}% divisibility-of-the-composite-lift
\chapters@set{20261126T030600}{36}% theorem-a-short-forms
\chapters@set{20261126T030700}{37}% the-determinant-of-a-chain
\chapters@set{20261126T030800}{38}% the-cost-of-one-lift
\chapters@set{20261126T030900}{39}% denominators-of-projections-to-lifts
\chapters@set{20261126T031000}{40}% twists-with-few-congruences
\chapters@set{20261126T031100}{41}% theorem-c-and-the-main-inequality
\providecommand\chapnumber[1]{%
  \@ifundefined{chapters@#1}%
    {\PackageError{chapters}{Sheet #1 is not listed in the index}%
      {List it in 20261125T012000--grh-workbook__guide.org\MessageBreak
       and run scripts/make-chapters.py.}}%
    {\edef\thechap{\@nameuse{chapters@#1}}}}
\makeatother
\chapnumber{<sheet ID>}
\makeatletter
\providecommand\chapters@set[2]{\@namedef{chapters@#1}{#2}}
\chapters@set{20261125T001000}{1}% zeros-to-ramp-sums
\chapters@set{20261125T001100}{2}% inert-primes-and-estimate-b
\chapters@set{20261125T005000}{3}% galois-fourier-and-averaging
\chapters@set{20261125T005500}{4}% columns-from-a-sphere
\chapters@set{20261125T000000}{5}% slope-method
\chapters@set{20261126T030100}{6}% a-twist-for-each-prime
\chapters@set{20261125T003000}{7}% the-chain-and-the-square-root
\chapters@set{20261125T010000}{8}% rank-lemma-and-zero-estimate
\chapters@set{20261125T010500}{9}% poincare-series-and-near-orthogonality
\chapters@set{20261124T223000}{10}% elliptic-curves-for-the-proof
\chapters@set{20261125T009500}{11}% fock-model-and-theta-intertwiners
\chapters@set{20261124T223100}{12}% complex-multiplication-and-inert-reduction
\chapters@set{20261125T008000}{13}% one-cm-determinant
\chapters@set{20261124T214500}{14}% supersingular-theta-divisibility
\chapters@set{20261124T214600}{15}% the-isotropic-plane
\chapters@set{20261125T013700}{16}% pairings-on-a-formal-group-of-height-two
\chapters@set{20261125T013500}{17}% hecke-nondegeneracy-on-the-tree
\chapters@set{20261126T030200}{18}% the-bound-is-attained
\chapters@set{20261125T009000}{19}% kronecker-limit-and-chowla-selberg
\chapters@set{20261125T002000}{20}% definite-theta-lifts
\chapters@set{20261125T015500}{21}% siegel-weil-in-the-smallest-case
\chapters@set{20261125T002100}{22}% definite-theta-lifts-and-norms
\chapters@set{20261125T016500}{23}% hecke-operators-and-theta-lifts
\chapters@set{20261125T016700}{24}% divisibility-along-the-chain
\chapters@set{20261125T004000}{25}% denominators-and-sha
\chapters@set{20261125T004100}{26}% tate-twists-and-sha
\chapters@set{20261126T020430}{27}% the-minus-part-by-gauss-sums
\chapters@set{20261125T014100}{28}% congruence-modules
\chapters@set{20261125T014500}{29}% eisenstein-congruences-and-selmer-classes
\chapters@set{20261125T014700}{30}% pseudocharacters-and-the-length-bound
\chapters@set{20261125T011000}{31}% galois-representations-black-box
\chapters@set{20261126T030300}{32}% the-data
\chapters@set{20261126T020210}{33}% bessel-poincare-series-and-their-gram-operators
\chapters@set{20261126T030400}{34}% the-seesaw-identity
\chapters@set{20261126T030500}{35}% divisibility-of-the-composite-lift
\chapters@set{20261126T030600}{36}% theorem-a-short-forms
\chapters@set{20261126T030700}{37}% the-determinant-of-a-chain
\chapters@set{20261126T030800}{38}% the-cost-of-one-lift
\chapters@set{20261126T030900}{39}% denominators-of-projections-to-lifts
\chapters@set{20261126T031000}{40}% twists-with-few-congruences
\chapters@set{20261126T031100}{41}% theorem-c-and-the-main-inequality
\providecommand\chapnumber[1]{%
  \@ifundefined{chapters@#1}%
    {\PackageError{chapters}{Sheet #1 is not listed in the index}%
      {List it in 20261125T012000--grh-workbook__guide.org\MessageBreak
       and run scripts/make-chapters.py.}}%
    {\edef\thechap{\@nameuse{chapters@#1}}}}
\makeatother
\chapnumber{20261125T001000}
\renewcommand{\thesection}{\thechap.\arabic{section}}

\theoremstyle{definition}
\newtheorem{problem}{Problem}[section]
\theoremstyle{plain}
\newtheorem{theorem}{Theorem}[section]

\newcommand{\Z}{\mathbb{Z}}
\newcommand{\Q}{\mathbb{Q}}
\newcommand{\C}{\mathbb{C}}
\newcommand{\R}{\mathbb{R}}
\renewcommand{\Re}{\operatorname{Re}}
\renewcommand{\Im}{\operatorname{Im}}
\newcommand{\one}{\mathbf{1}}

\setlength{\emergencystretch}{2em}

\title{From zeros to a one-sided bound}
\author{}
\date{November 26, 2026}

\begin{document}
\maketitle
\refstepcounter{chap}\label{chap}


\paragraph{Prerequisites.}
The zeta function and Dirichlet $L$-functions as in \cite{Dav}: orthogonality of characters and $L(1, \chi) \ne 0$ for nonprincipal $\chi$ (Ch.~4), the class number formula (Ch.~6), the Hadamard product and the constant $B$ (Ch.~12), nonvanishing on $\Re s = 1$ (Ch.~13 and~14), the explicit formula for Chebyshev's function $\psi(x)$ and its contour shift (Ch.~17), their analogues for $L(s, \chi)$ (Ch.~12 and~19), and the prime number theorem in progressions to a fixed modulus (Ch.~20).  Chebyshev's bound $\psi(x) \ll x$ \cite[\S2.2]{MV}.  Mellin transforms.  Landau's theorem is stated in \S\ref{sec:landau}.  Section~\ref{sec:dirichlet} uses the auxiliary-prime theorem stated as Theorem~\ref{ib-thm:choice} in Chapter~\ref{ib-chap}.  Chapter~\ref{ip-chap} is used only for motivation, in \S\ref{sec:smoothing}. % check-numbers: ok (Davenport): Ch.~4; Ch.~6; Ch.~12; Ch.~13 and~14; Ch.~17; Ch.~12 and~19; Ch.~20

\medskip
A comparison of divisibility with size, as in Chebyshev's method, gives upper bounds for weighted sums over primes.  Can an upper bound for a sum over primes prove the Riemann hypothesis?  This chapter shows that it can, for a suitable weight.  After the zeta function, Theorem~\ref{thm:dirichlet} gives the criterion for Dirichlet $L$-functions in the form used in the proof.  The upper bound itself is not treated here.

\section{The ramp sum}\label{sec:ramp}

Write $v_+ := \max(v, 0)$ for real $v$, and let $\Lambda$ be the von Mangoldt function.  The \emph{ramp} at $x$ is the function $n \mapsto (x/n - 1)_+$: it is large for small $n$, decreases, and vanishes for $n \ge x$.  Let
\begin{equation}\label{eq:G}
  S(x) := \sum_{n} \Lambda(n) \Bigl(\frac{x}{n} - 1\Bigr)_+ ,
  \qquad
  G(x) := S(x) - x \log x + (1 + \gamma) x ,
\end{equation}
where $\gamma$ is Euler's constant, and let $\psi(x) := \sum_{n \le x} \Lambda(n)$ be Chebyshev's function.  In \S\ref{sec:dirichlet}, $\psi$ denotes a Dirichlet character.  The two subtracted terms are the main term of $S$, as the problems show.

\begin{problem}\label{pr:ramp-mellin}
For $n \ge 1$ and $\Re s > 1$,
\(
  \int_1^\infty (x/n - 1)_+ \, x^{-s-1} \, d x = n^{-s}/\bigl(s(s-1)\bigr)
\).
\end{problem}

\begin{problem}\label{pr:S-integral}
$S(x) = x \int_1^x \psi(t) \, t^{-2} \, d t$.  Hence $G(x) \ll x \log x$ for $x \ge 2$.
\end{problem}

\begin{problem}\label{pr:F}
For $\Re s > 1$,
\begin{equation}\label{eq:F}
  F(s) := \int_1^\infty G(x)\, x^{-s-1} \, d x
  = -\frac{\zeta'}{\zeta}(s) \, \frac{1}{s(s-1)} - \frac{1}{(s-1)^2} + \frac{1 + \gamma}{s - 1} .
\end{equation}
\end{problem}

\begin{problem}\label{pr:F-poles}
$F$ is meromorphic in $\C$ and regular at $s = 1$.  Its poles in $\Re s > 0$ are the zeros $\rho$ of $\zeta$ in the critical strip $0 \le \Re s \le 1$, and a zero of multiplicity $m$ gives residue $-m/\bigl(\rho(\rho - 1)\bigr)$.
\end{problem}

\begin{problem}\label{pr:residue-zero}
The residue of $F(s) x^s$ at $s = 0$ is $\log 2\pi$.
\end{problem}

Problems~\ref{pr:F}--\ref{pr:residue-zero} and a contour shift as in \cite[Ch.~17]{Dav}, applied to $-\frac{\zeta'}{\zeta}(s)\, x^s/(s(s-1))$, the part of $F(s) x^s$ that comes from $S$, where now every integral converges absolutely because of the factor $1/(s(s-1))$, give the following.

\begin{theorem}\label{thm:explicit}
For $x > 1$,
\[
  G(x) = \log 2\pi - \sum_\rho \frac{x^\rho}{\rho(\rho - 1)} + O(x^{-2}) ,
\]
where $\rho$ runs over the zeros of $\zeta$ in the critical strip with multiplicity, and the sum converges absolutely.
\end{theorem}

\begin{problem}\label{pr:sum-rho}
$\sum_\rho 1/\bigl(\rho(1 - \rho)\bigr) = 2 + \gamma - \log 4\pi = 0.0461\ldots$, with multiplicity.
\end{problem}

\begin{problem}\label{pr:rh-bound}
If the Riemann hypothesis holds, then $\lvert G(x) - \log 2\pi \rvert \le (2 + \gamma - \log 4\pi)\sqrt{x} + O(x^{-2})$.
\end{problem}

For instance $G(10^4) \approx 1.749$, $G(10^5) \approx 3.671$, $G(10^6) \approx 11.040$ and $G(2 \cdot 10^6) \approx -12.268$, to be compared with $\log 2\pi \approx 1.838$ and $0.0462 \sqrt{10^6} = 46.2$.

\begin{problem}\label{pr:one-pair}
Let $\rho = \beta + i\gamma_0$ be a nonreal zero of $\zeta$, and let $c := -1/\bigl(\rho(\rho-1)\bigr)$.  The joint term $c x^\rho + \bar c x^{\bar\rho}$ of $\rho$ and $\bar\rho$ in Theorem~\ref{thm:explicit} exceeds $\lvert c \rvert x^\beta$ for all $x$ in a set whose image under $\log$ has density $1/3$ in $[0, \infty)$, and it is less than $-\lvert c \rvert x^\beta$ on another such set.
\end{problem}

If $\beta > \frac12$, the contribution of this conjugate pair exceeds $\sqrt x$ along one sequence tending to infinity and falls below $-\sqrt x$ along another.  A comparison of divisibility with size sees only one sign.  The next section shows that one is enough.

\section{One side suffices}\label{sec:landau}

We use the following measurable version of Landau's theorem \cite{Lan}.  Montgomery and Vaughan \cite[Lemma~15.1]{MV} assume local Riemann integrability and use $x^{-s}$; their proof applies to the Lebesgue integral below after replacing $s$ by $s+1$.  Let $A \colon [1, \infty) \to \R$ be measurable and bounded on compact sets, with $A(x) \ge 0$ for $x \ge X_0$.  Let $\sigma_c$ be the infimum of the real $\sigma$ with $\int_1^\infty \lvert A(x) \rvert x^{-\sigma-1} \, d x < \infty$, and suppose that $-\infty < \sigma_c < \infty$.  Then $H(s) := \int_1^\infty A(x) x^{-s-1} \, d x$ is holomorphic in $\Re s > \sigma_c$, and it has no holomorphic continuation to a neighborhood of $\sigma_c$.

\begin{problem}\label{pr:below}
In this situation $H(\sigma)$ is bounded below for real $\sigma \in (\sigma_c, \sigma_c + 1]$.
\end{problem}

Fix $0 < \theta < 1$, suppose that $G(x) \le x^\theta$ for all large $x$, and let $A(x) := x^\theta - G(x)$.

\begin{problem}\label{pr:H}
$-\infty < \sigma_c \le 1$, and $H(s) = 1/(s - \theta) - F(s)$ for $\Re s > 1$.
\end{problem}

\begin{problem}\label{pr:zeta-negative}
$\zeta(\sigma) < 0$ for $0 < \sigma < 1$.
\end{problem}

\begin{problem}\label{pr:real-zero}
If $\sigma_c > \theta$, then $\sigma_c$ is a real zero of $\zeta$, and $H(\sigma) \to -\infty$ as $\sigma \downarrow \sigma_c$.  Hence $\sigma_c \le \theta$.
\end{problem}

By Problem~\ref{pr:zeta-negative} $\zeta$ has no real zero in $(0,1)$, but Problem~\ref{pr:real-zero} does not use this, and in \S\ref{sec:dirichlet} the same argument will be applied where real zeros cannot be excluded in advance.  By Problem~\ref{pr:real-zero}, $F = 1/(s-\theta) - H$ is holomorphic in $\Re s > \theta$, and by Problem~\ref{pr:F-poles} this excludes zeros there.

\begin{theorem}\label{thm:one-sided}
Let $0 < \theta < 1$.  If $G(x) \le x^\theta$ for all large $x$, then $\zeta$ has no zeros in $\Re s > \theta$.
\end{theorem}

\begin{problem}\label{pr:rh-from}
If for every $\varepsilon > 0$ one has $G(x) \le x^{1/2 + \varepsilon}$ for all large $x$, then the Riemann hypothesis holds.
\end{problem}

The transform $F$ continues meromorphically, with at most simple poles in $\Re s > 0$.  This, the sign of its residues at real points to the right of $\theta$, and the eventual upper bound for $G$ are what the argument of Problem~\ref{pr:real-zero} uses.  In the variable $v = \log x$, with $E(v) := e^{-v} G(e^v)$, the transform $F(s)$ becomes $\int_0^\infty E(v) e^{(1-s)v} \, d v$, and the bound $G(x) \le x^\theta$ becomes $E(v) \le e^{(\theta - 1)v}$.  In this form the argument proves the following.

\begin{theorem}\label{thm:landau-polar}
Let $E \colon \R \to \R$ and $r \colon \C \to \C$.  Suppose that $E$ is Borel measurable, and that there are $C \in \R$ and a meromorphic function $\Phi$ on $\C$ such that $\lvert E(v) \rvert \le C e^v$ for all $v \ge 0$, such that $\Phi(s) = \int_0^\infty E(v) e^{(1-s)v} \, d v$ for all $s$ with $\Re s > 2$, and such that $(z - \rho)\Phi(z) \to r(\rho)$ as $z \to \rho$, $z \ne \rho$, for every $\rho$ with $\Re \rho > 0$.  Let $\beta > 0$.  Suppose that $\Re r(x) \ge 0$ for every real $x > \beta$, and that there is $V$ with $E(v) \le e^{(\beta - 1)v}$ for all $v \ge V$.  Then $r(\rho) = 0$ for every $\rho \in \C$ with $\Re \rho > \beta$.
\end{theorem}

The hypothesis on $\Phi$ says that its poles in $\Re s > 0$ are simple, with residues $r$; the conclusion says that it has none in $\Re s > \beta$.  Only a sign is assumed at real points, and only there.

\section{Smoothing, and the triangle weight}\label{sec:smoothing}

Why the ramp, and not another weight?  Let $\phi \colon (0, \infty) \to \R$ be measurable, vanish on $(1, \infty)$, be bounded on compact subsets of $(0, 1]$, and satisfy $\phi(t) = O(1/t)$ as $t \to 0$.  Let
\[
  \psi_\phi(x) := \sum_n \Lambda(n) \phi(n/x) ,
  \qquad
  \tilde\phi(s) := \int_0^\infty \phi(t) t^{s-1} \, d t \quad (\Re s > 1) .
\]
We use $\tilde\phi$ on its half-plane of absolute convergence and, for the three weights below, also for the meromorphic continuation given by their rational formulas.  The sharp cutoff $\phi = \one_{(0,1]}$ gives $\psi_\phi = \psi$ and $\tilde\phi(s) = 1/s$.  The ramp is $\phi(t) = (1/t - 1)_+$, with $\tilde\phi(s) = 1/(s(s-1))$ by Problem~\ref{pr:ramp-mellin}.  The \emph{triangle weight} is $\phi(t) = (1 - t)_+$, with $\psi_\phi(x) = x^{-1}\int_0^x \psi(t) \, d t$.

\begin{problem}\label{pr:smooth-mellin}
For $\Re s > 1$, $\int_0^\infty \psi_\phi(x) x^{-s-1} \, d x = -\frac{\zeta'}{\zeta}(s) \, \tilde\phi(s)$.
\end{problem}

\begin{problem}\label{pr:triangle-mellin}
For the triangle weight $\tilde\phi(s) = 1/(s(s+1))$ for $\Re s > 0$.  If $\phi_1(t) := \phi(t)/t$, then $\tilde\phi_1(s) = \tilde\phi(s - 1)$ wherever either integral converges absolutely; for the triangle weight $\phi_1$ is the ramp.
\end{problem}

The following theorem is proved as Theorem~\ref{thm:explicit}: by Problem~\ref{pr:smooth-mellin} with $\tilde\phi(s) = 1/(s(s+1))$, shift the line of integration to $\Re s = -3/2$ as in \cite[Ch.~17]{Dav}, along horizontal segments at heights $T$ where $\zeta'/\zeta \ll \log^2 T$.  The factor $1/(s(s+1))$ makes every integral converge absolutely, and the poles of $-\frac{\zeta'}{\zeta}(s)\, x^s/(s(s+1))$ to the right of $-3/2$ are at $1$, $0$, $-1$ and the zeros $\rho$.

\begin{theorem}\label{thm:triangle}
For $x > 1$,
\[
  \sum_{n \le x} \Lambda(n) \Bigl(1 - \frac{n}{x}\Bigr) = \frac{x}{2} - \log 2\pi - \sum_\rho \frac{x^\rho}{\rho(\rho + 1)} + O\Bigl(\frac1x\Bigr) .
\]
\end{theorem}

As a check, the left side minus $(x/2 - \log 2\pi)$ is approximately $0.106$, $7.975$ and $-12.115$ at $x = 10^4$, $10^6$ and $2 \cdot 10^6$.  For the ramp and the triangle weight $\lvert \tilde\phi(\rho) \rvert \ll \lvert \rho \rvert^{-2}$, so the sums over zeros converge absolutely.  For the sharp cutoff they do not, since $\sum_\rho 1/\lvert \rho \rvert$ diverges.

\begin{problem}\label{pr:triangle-one-sided}
Let $0 < \theta < 1$.  If $\sum_{n \le x} \Lambda(n)(1 - n/x) \le x/2 + x^\theta$ for all large $x$, then $\zeta$ has no zeros in $\Re s > \theta$.
\end{problem}

We return to real zeros in \S\ref{sec:dirichlet}.

The ramp itself is forced by the arithmetic.  Theorem~\ref{ip-thm:main} gives, at primes $p \equiv 3 \pmod 8$, lower bounds for $p$-adic valuations of the form $\sum_{i \text{ odd}} (\lfloor a/p^i \rfloor - w)_+$.  Up to an error of $1$ for each odd $i$ with $p^i \le a$, and with $a = wx$, the summand is $(a/p^i - w)_+ = w(x/p^i - 1)_+$, which is $w$ times the ramp at $x$.  Bounds of this form at every inert prime, summed with weights $\log p$, give the ramp sum over odd powers of primes inert in an imaginary quadratic field; Chapter~\ref{ib-chap} says what this sum counts.

\section{Dirichlet \texorpdfstring{$L$}{L}-functions}\label{sec:dirichlet}

What replaces $G$ for the $L$-functions of the characters modulo $q$?  The arithmetic bounds sums over primes that do not split in an imaginary quadratic field.  For a prime $d \equiv 3 \pmod 4$, the Legendre symbol $\chi_d(n) = (n/d)$ is the character of $K = \Q(\sqrt{-d})$, and at a prime $p$ the weight $1 - \chi_d(p)$ is $2$ if $p$ is inert, $1$ if $p = d$, and $0$ if $p$ splits.  This weight brings in the twisted $L$-functions $L(s, \psi\chi_d)$, whose orders of vanishing are subtracted from the untwisted ones in the residues below.  The prime $d$ will be chosen so that they do not vanish where it matters.

For a primitive nonprincipal Dirichlet character $\chi$ let $c_\chi := \frac{L'}{L}(1, \chi)$ and
\begin{equation}\label{eq:Gchi}
  G_\chi(x) := \sum_{n \le x} \Lambda(n) \chi(n) \Bigl(\frac{x}{n} - 1\Bigr) + c_\chi x .
\end{equation}
Let $d$ be an odd prime, $\chi_d(n) := (n/d)$ the Legendre symbol, and $q \ge 1$ an integer prime to $d$.  Let $\varphi$ be Euler's function.  For $n \ge 1$ let
\[
  \lambda(n) := \Lambda(n)\bigl(1 - \chi_d(n)\bigr) \ \text{if } n \equiv 1 \ (\mathrm{mod}\ q),
  \qquad \lambda(n) := 0 \ \text{otherwise}.
\]
The analytic definitions below apply to every odd prime $d$.  For a Dirichlet character $\psi$ modulo $q$, let $L(s, \psi)$ be the meromorphic continuation of $\sum_{n \ge 1} \psi(n) n^{-s}$ from $\Re s > 1$; its Euler product there runs over the primes $p \nmid q$.  Let $E_\psi := L(\,\cdot\,, \psi)$ if $\psi$ is not principal, and if $\psi$ is principal let $E_\psi$ be the entire continuation of $(s - 1)L(s, \psi)$, with $E_\psi(1) = \prod_{p \mid q}(1 - p^{-1})$.  Let $m_\psi(\rho) \ge 0$ be the order of vanishing of $E_\psi$ at $\rho$.  The product $\psi\chi_d$ is a character modulo $qd$; it is not principal, and we write $E_{\psi\chi_d} := L(\,\cdot\,, \psi\chi_d)$ and $m_{\psi\chi_d}$ for its order.  Let
\[
  G_{d,q}(s) := \frac{1}{\varphi(q)} \Bigl( \sum_{\psi} \Bigl( \frac{E_{\psi\chi_d}'}{E_{\psi\chi_d}}(s) - \frac{E_\psi'}{E_\psi}(s) \Bigr) - 1 \Bigr),
\]
the sum over the characters $\psi$ modulo $q$.  The $E$'s do not vanish on $\Re s \ge 1$, so $G_{d,q}$ is holomorphic there.  Let $c_{d,q} := \Re G_{d,q}(1)$, and for real $x > 0$ let
\begin{equation}\label{eq:R}
  R_{d,q}(x) := \sum_{1 \le n \le x} \lambda(n) \Bigl( \frac1n - \frac1x \Bigr) - \frac{\log x}{\varphi(q)} - c_{d,q} .
\end{equation}
The simplest case is $q = 1$, where $\lambda(n) = \Lambda(n)(1 - \chi_d(n))$ for all $n$ and $E_{\one}(s) = (s - 1)\zeta(s)$ for $s \ne 1$, with $E_{\one}(1) = 1$.

\begin{problem}\label{pr:coef}
$\lambda(n) = \varphi(q)^{-1} \sum_\psi \bigl(\psi(n) - \psi(n)\chi_d(n)\bigr) \Lambda(n)$, and $0 \le \lambda(n) \le 2\Lambda(n)$.  For $q = 1$ and $d = 7$ the nonzero values for $n \le 11$ are $\lambda(3) = 2\log 3$, $\lambda(5) = 2\log 5$ and $\lambda(7) = \log 7$.
\end{problem}

\begin{problem}\label{pr:dirichlet-series}
For $\Re s > 1$,
\(
  \sum_{n \ge 1} \lambda(n) n^{-s} = s/\bigl(\varphi(q)(s-1)\bigr) + G_{d,q}(s)
\).
\end{problem}

\begin{problem}\label{pr:R-transform}
For $\Re s > 1$,
\(
  \int_1^\infty R_{d,q}(y)\, y^{-s} \, d y = \bigl(G_{d,q}(s)/s - c_{d,q}\bigr)/(s - 1)
\).
\end{problem}

\begin{problem}\label{pr:product}
(i) The function $\Phi(s) := \bigl(G_{d,q}(s)/s - c_{d,q}\bigr)/(s - 1)$ is meromorphic on $\C$ and regular at $s = 1$.  At every $\rho \ne 0, 1$ it has at most a simple pole, with residue
\[
  \frac{1}{\varphi(q)\, \rho(1 - \rho)} \sum_\psi \bigl( m_\psi(\rho) - m_{\psi\chi_d}(\rho) \bigr) .
\]
(ii) If no $L(s, \psi\chi_d)$ vanishes at a real $x \in (0, 1)$, the residue at $x$ is $\ge 0$, and it is $> 0$ if some $L(s, \psi)$ vanishes there.  A zero of order $m$ of one $L(s, \psi\chi_d)$ at a point $\rho \ne 0, 1$ where no other $E$ vanishes gives the residue $-m/\bigl(\varphi(q)\rho(1-\rho)\bigr)$.
\end{problem}

At $\rho \ne 0, 1$ the residue depends on the difference between the total untwisted and twisted orders.  A common zero can therefore reduce or cancel the contribution of an untwisted factor, and at a real point $0 < \rho < 1$ a twisted zero alone gives a negative residue.  Neither happens where no twist vanishes.

\begin{problem}\label{pr:c-value}
$c_{d,1} = -1 - \gamma + \frac{L'}{L}(1, \chi_d)$.  For $d = 7$, $L(1, \chi_7) = \pi/\sqrt 7 = 1.1874104\ldots$; with $\frac{L'}{L}(1, \chi_7) \approx 0.0156357$ this gives $c_{7,1} \approx -1.56158$.
\end{problem}

\begin{problem}\label{pr:G-difference}
For $y \ge 1$, $y R_{d,1}(y) = G(y) - G_{\chi_d}(y)$, with $G$ as in \eqref{eq:G} and $G_{\chi_d}$ as in \eqref{eq:Gchi}.
\end{problem}

For instance $R_{7,1}(y) \approx 6.27 \cdot 10^{-4}$, $1.05 \cdot 10^{-4}$ and $-7.71 \cdot 10^{-5}$ at $y = 10^4$, $10^5$ and $10^6$.

\begin{problem}\label{pr:R-limit}
(i) $\sum_{n \le y} \lambda(n)(1/n - 1/y) - \log y/\varphi(q)$ has a limit as $y \to \infty$.  (ii) The limit is $c_{d,q}$; that is, $R_{d,q}(y) \to 0$.
\end{problem}

So $c_{d,q}$ is the constant that makes the remainder small.  The proof gives more: for every odd prime $d$, every $q \ge 1$ prime to $d$ and every integer $B \ge 0$, $\lvert R_{d,q}(y) \rvert \le (\log y)^{-B}$ for all large real $y$.  It does this by a Tauberian theorem applied to $v \mapsto R_{d,q}(e^v)$ and the transform identity of Problem~\ref{pr:R-transform} for $\Re s > 2$.

Now fix a point $\rho$.  The next problem is the analogue of Theorem~\ref{thm:one-sided} for the $L(s, \psi)$, at one prime $d$ whose twists do not vanish near the real axis up to the height of $\rho$.

\begin{problem}\label{pr:fixed-d}
Let $d \equiv 7 \pmod 8$ be prime, $q \ge 1$ prime to $d$, $\beta > 0$, and $\rho \ne 1$ with $\Re \rho > \beta$.  Suppose that $L(s, \psi\chi_d) \ne 0$ for every $\psi$ modulo $q$ and every $s$ with $\Re s \ge \beta$ and $\lvert \Im s \rvert \le \lvert \Im \rho \rvert$, and that $R_{d,q}(y) \le y^{\beta - 1}$ for all large real $y$.  Then $L(\rho, \psi) \ne 0$ for every $\psi$ modulo $q$.
\end{problem}

Problem~\ref{pr:fixed-d} assumes no separate exclusion of real zeros of the untwisted $L$-functions.  For $\beta > 1/2$, Theorem~\ref{ib-thm:choice} of Chapter~\ref{ib-chap} supplies the required twist nonvanishing for some prime $d$ in any residue class satisfying its hypotheses.  For a Dirichlet character $\xi$ modulo $N$ and an odd prime $d$ not dividing $N$, write $\xi\chi_d$ for their product, a character modulo $Nd$.  That theorem says: let $N \ge 1$, $\beta > 1/2$ and $t_0 \in \R$, and let $a$ be a unit modulo $N$ such that every prime $p \equiv a \pmod N$ satisfies $p \equiv 7 \pmod 8$ and $\gcd(N, p) = 1$; then there is a prime $d \equiv a \pmod N$ such that $L(s, \xi\chi_d) \ne 0$ for every character $\xi$ modulo $N$ and every $s$ with $\Re s \ge \beta$ and $\lvert \Im s \rvert \le t_0$.  It says nothing about $R_{d,q}$.

Let $q_0$ be a positive multiple of $8$, and let $d_0 \equiv 7 \pmod 8$ be prime to $q_0$.  Every prime $d \equiv d_0 \pmod{q_0}$ is then $\equiv 7 \pmod 8$ and prime to $q_0$.  The next theorem assumes the bound for $R_{d,q_0}$ for every prime $d$ in this class.  Given $\rho$, Theorem~\ref{ib-thm:choice} with $N = q_0$, $a = d_0$ and $t_0 = \lvert \Im \rho \rvert$ chooses $d$ in the class, and Problem~\ref{pr:fixed-d} applies to it.

\begin{theorem}\label{thm:dirichlet}
Let $q_0 \ge 1$ with $8 \mid q_0$, let $d_0 \ge 0$ be an integer with $d_0 \equiv 7 \pmod 8$ and $\gcd(d_0, q_0) = 1$, and let $\beta > 1/2$.  Suppose that for every prime $d \equiv d_0 \pmod{q_0}$ there is $Y \in \R$ such that $R_{d,q_0}(x) \le x^{\beta - 1}$ for every real $x \ge Y$.  Then $L(\rho, \psi) \ne 0$ for every character $\psi$ modulo $q_0$ and every $\rho \ne 1$ with $\Re \rho > \beta$.
\end{theorem}

The hypothesis concerns every prime of one class modulo $q_0$, and the point $Y$ from which the bound holds may depend on $d$.  For $q = 1$, by Problem~\ref{pr:G-difference}, a bound $R_{d,1}(x) \le x^{\theta - 1}$ reads $G(x) - G_{\chi_d}(x) \le x^\theta$; compare Theorem~\ref{thm:one-sided}.

\begin{problem}\label{pr:rh-dirichlet}
Let $M \ge 1$.  If for every $\beta > 1/2$ the hypothesis of Theorem~\ref{thm:dirichlet} holds for $q_0 = 8M$ and $d_0 = 8M - 1$, then $L(s, \chi) \ne 0$ for every character $\chi$ modulo $M$ and every $s \ne 1$ with $\Re s > 1/2$.
\end{problem}

\paragraph{In the Lean formalization.}
{\sloppy In the Lean proof \cite{GL} (readable edition), Theorem~\ref{thm:dirichlet} is \texttt{New.LFunction\_ne\_zero\_of\_primeRemainder\_le} in \texttt{New/Reduction/Zeros/Criterion.lean}; for $x > 0$, $R_{d,q}(x)$ is \texttt{New.primeRemainder} at $d, q, x$ in \texttt{New/Reduction/Zeros/PrimeError.lean}.  Theorem~\ref{thm:landau-polar} is \texttt{New.IsPolar.landau} in \texttt{Lib/Analytic/Landau.lean}, with the definition \texttt{New.IsPolar} written out.  The bound $\lvert R_{d,q}(y) \rvert \le (\log y)^{-B}$ is \texttt{New.eventually\_abs\_primeRemainder\_le} in \texttt{New/Reduction/Zeros/PrimeNumberTheorem.lean}.\par}

\section*{Hints}

\subsection*{\S\ref{sec:ramp}}
\begin{description}
\item[\ref{pr:S-integral}] Integrate each step of $\psi$.  For the bound use Chebyshev's estimate.
\item[\ref{pr:F}] Use Problem~\ref{pr:ramp-mellin} termwise, and compute the transforms of $x \log x$ and $x$.
\item[\ref{pr:F-poles}] Use $-\zeta'/\zeta(s) = 1/(s-1) - \gamma + O(s - 1)$ and expand $1/(s(s-1))$ at $s = 1$.
\item[\ref{pr:residue-zero}] The functional equation gives $\zeta'/\zeta(0) = \log 2\pi$.
\item[\ref{pr:sum-rho}] Write $1/(\rho(1-\rho)) = 1/\rho + 1/(1 - \rho)$ and use the value of $B$ in \cite[Ch.~12]{Dav}.
\end{description}

\subsection*{\S\ref{sec:landau}}
\begin{description}
\item[\ref{pr:below}] Split the integral at $X_0$.
\item[\ref{pr:H}] Use Problem~\ref{pr:S-integral} for $\sigma_c \le 1$ and Problem~\ref{pr:F} for the formula for $H$.  For $\sigma_c > -\infty$ use Problem~\ref{pr:F-poles}: $\zeta$ has infinitely many zeros in the critical strip \cite[Ch.~12]{Dav}, so it has nonreal ones.
\item[\ref{pr:zeta-negative}] Consider $(1 - 2^{1-\sigma})\zeta(\sigma) = \sum_{n \ge 1} (-1)^{n-1} n^{-\sigma}$.
\item[\ref{pr:real-zero}] By Landau's theorem $\sigma_c$ is a singularity of $H$.  Of the real points in $(\theta, 1]$, only real zeros of $\zeta$ can be singular (Problem~\ref{pr:F-poles}).  Find the sign of the residue of $F$ at a real zero $\beta \in (0,1)$, and compare with Problem~\ref{pr:below}.
\item[\ref{pr:rh-from}] Use Theorem~\ref{thm:one-sided} and the symmetry $\rho \mapsto 1 - \rho$.
\end{description}

\subsection*{\S\ref{sec:smoothing}}
\begin{description}
\item[\ref{pr:smooth-mellin}] Substitute $x = n/t$ in each term.
\item[\ref{pr:triangle-one-sided}] Repeat \S\ref{sec:landau} with $\tilde\phi(s) = 1/(s(s+1))$, including the argument for $\sigma_c > -\infty$ in Problem~\ref{pr:H}, and compare the sign of the residue at a real zero with that in Problem~\ref{pr:real-zero}; conclude with Problem~\ref{pr:zeta-negative}.
\end{description}

\subsection*{\S\ref{sec:dirichlet}}
\begin{description}
\item[\ref{pr:coef}] Orthogonality of the characters modulo $q$; a number $n \equiv 1 \pmod q$ is prime to $q$.
\item[\ref{pr:dirichlet-series}] Problem~\ref{pr:coef} and $-\frac{L'}{L}(s, \chi) = \sum_n \chi(n)\Lambda(n) n^{-s}$; for the principal $\psi$, $E_\psi'/E_\psi$ and $L'/L$ differ by $1/(s-1)$.
\item[\ref{pr:R-transform}] As Problem~\ref{pr:F}: $(1/n - 1/y)_+ y^{-s} = (y/n - 1)_+ y^{-s-1}$, so Problem~\ref{pr:ramp-mellin} applies; then Problem~\ref{pr:dirichlet-series}.
\item[\ref{pr:product}] (i) A zero of order $m$ of an entire function at $\rho$ is a simple pole of its logarithmic derivative, with residue $m$.  At $s = 1$: the characters modulo $q$ are closed under complex conjugation, so $G_{d,q}(1)$ is real.
\item[\ref{pr:c-value}] For $q = 1$ the untwisted factor is $(s-1)\zeta(s) = 1 + \gamma(s - 1) + \cdots$.  For $d = 7$ use the class number formula.
\item[\ref{pr:G-difference}] Problem~\ref{pr:c-value}; for $n \le y$ the ramp is $y(1/n - 1/y)$.
\item[\ref{pr:R-limit}] (i) Partial summation, as in Problem~\ref{pr:S-integral}, with the prime number theorem with error term for the characters modulo $qd$ \cite[Ch.~20]{Dav}.  (ii) If $R_{d,q} \to \ell$, compare the behaviour of the two sides of Problem~\ref{pr:R-transform} as $s \downarrow 1$, using Problem~\ref{pr:product}.
\item[\ref{pr:fixed-d}] Theorem~\ref{thm:landau-polar} for $E(v) := R_{d,q}(e^v)$, with Problems~\ref{pr:R-transform} and~\ref{pr:product}.  Check the sign of the residue at real $x$ in $(\beta, 1)$, at $1$ and in $(1, \infty)$ separately.
\item[\ref{pr:rh-dirichlet}] A character modulo $M$ induces one modulo $8M$; the two $L$-functions differ by finitely many Euler factors.
\end{description}

\begin{thebibliography}{Alp}

\bibitem[Dav]{Dav}
H.~Davenport, \emph{Multiplicative Number Theory}, 3rd ed., revised by H.~L.~Montgomery, Grad. Texts in Math. \textbf{74}, Springer, 2000.

\bibitem[GL]{GL} L.~Sernau, \emph{GRH for Dirichlet \(L\)-functions, in Lean~4}, a shorter rewriting of the Lean proof, text by Claude, \url{https://github.com/lsernau/grh-lean}, commit \texttt{72fead1}, readable edition \texttt{grh/}.

\bibitem[Lan]{Lan}
E.~Landau, \emph{\"Uber einen Satz von Tschebyschef}, Math. Ann. \textbf{61} (1905), 527--550.

\bibitem[MV]{MV}
H.~L.~Montgomery and R.~C.~Vaughan, \emph{Multiplicative Number Theory~I. Classical Theory}, Cambridge Stud. Adv. Math. \textbf{97}, Cambridge University Press, 2007.

\end{thebibliography}

\end{document}
